% This file should be rename with the speaker's last name. (e.g : smith)

% Write the author names in the order you wish them to appear
% and underline the one who will give the talk.

% Only the speaker's email address is required.

\begin{abstract_online}{Interactive tutorials, an example on symmetric functions}{%
        \underline{Pauline Hubert}$^{1}$, M??lodie Lapointe$^{1}$ }{[hubert.pauline@courrier.uqam.ca]
        }{%
        $^1$D??partement de Math??matiques, Universit?? du Qu??bec ?? Montr??al, Montr??al, Canada}
    
Sage is a free open source computer algebra software. The project was started in 2005 by William Stein [2] as an open source alternative to mathematical systems such as Maple or Mathematica, and is based on python and Jupyter. Thanks to its hundreds of worldwide contributors, Sage now contains a large variety of libraries such as calculus, linear algebra, combinatorics, number theory, and it is used intensively in research and higher education. 

Current tutorials and documentation are often written by top specialists in their fields, because of this, it can be hard to access for newcomers. Thus, we wanted to build a tutorial which is both a mathematical introduction to the subject, and a tutorial on how to use the relevant tools in Sage. \\
%We assumed few previous mathematical knowledge of the subject by the user, and very basic knowledge of Sage itself. \\
In our case, we have been mostly interested in the symmetric function tools. The classical mathematical reference here is [1]. Our goals were to improve and complete the pre-existent tutorials, to add an interactive dimension and and to show the mathematics behind and not only the Sage tools. The expected result would thus interlace class notes with an actual tutorial on how to use Sage to explore the notions considered.

In this presentation, we will use the example of this tutorial to present some interesting features of Sage and Jupyter. We will also talk about how interactive tutorials and notebooks may be turned into learning tools.
One of the key feature here the closeness between the mathematical development of the subject considered and the Sage programming style. \\


\textbf{Keywords} \newline{} Interactive tutorial, Sage, Symmetric functions

\textbf{References} \newline{}
% Model for references to books.
[1] \textsc{I. G. Macdonald}, \emph{Symmetric functions and Hall polynomials}. Oxford University Press, New York, 1995. \newline{}
%  Model for references to articles
[2] \textsc{W. Stein}, Sage: creating a viable free open source alternative to Magma, Maple, Mathematica, and MATLAB. \emph{London Math. Soc. Lecture Note Ser.} \textbf{403}, 230--238  (2013). \newline{}
\end{abstract_online}
    